\documentclass[exercice]{thedocument}%
\usepackage{xlop}%
\startdatakeys
    \source{}
    \BO{2026}
    \cycle{3}
    \competencesDuSocle{CA1}
    \objectifsApprentissage{C3A3ax}
    \motsCles{division euclidienne, quotient, reste}
\enddatakeys
\begin{document}%
    \!n1=112;%
    \!n2=6;%
    \begin{center}
        \opidiv[displayintermediary=all,voperation=top]{\<n1;}{\<n2;}
        \end{center}
        \vspace{.5cm}
    \begin{enumerate}
        \item Dans la division euclidienne ci-contre, quel est le reste ? Le quotient ?
        Le dividende ? Le diviseur ?
        \item Quelle égalité peut-on écrire ?
    \end{enumerate}
\end{document}
